\documentclass{article}
\usepackage{amsmath}
\usepackage{amssymb}
\usepackage{here}
\usepackage[all]{foreign}
\usepackage{graphicx}
\usepackage[colorlinks=true, allcolors=blue]{hyperref}
\title{Cover Letter for: \textit{Pointing Beyond Fitts: Strategy, Movement Time Distributions, Local-Global Tradeoffs, and Generative Models}}
\author{Julien Gori}
\begin{document}
\maketitle


Dear Dr. Pourang Irani,

\vspace{2\baselineskip}

Thank you for the detailed reviews, and for giving me the chance to make the manuscript better. I was happy to see that the paper did spark interest, despite the criticism. The first round of reviews has been very constructive and I believe the paper is now stronger, and I hope the reviewers will agree.

\vspace{\baselineskip}

I agree with many of the points raised in the reviews, and in particular with the high level comments regarding the overall structure of the paper and the narrative, which indeed needed to be reworked to make the research contributions more apparent.
I have edited many parts of the manuscript, in particular the first sections that introduce and motivate the problem, as well as the discussion section, but no section was left untouched. I have also removed content that was secondary, and moved yet other content to the supplementary materials. At the bottom of this letter, I provide point by point responses to all reviewer comments.

I hope you and external reviewers will find that I have addressed the concerns brought to light in this first round, and that the goal and contributions of the paper are now clear. 

\begin{flushright}
Sincerely,\\
Julien Gori
\end{flushright}



\section*{Answering reviewer comments}

\subsection*{Reviewer 1 comments}

\begin{enumerate}
    \item \textit{Structure/Rationale: There is a lot in this work, including the introduction of copulas, the fitting of (1) copulas, (2) EMG models, and (3) bivariate Gaussian models for different datasets and study protocols, the comparison between copulas, the comparison between local and global fits, and other more minor contributions. While all this is interesting to read, and the paper has a nice flow, the relationships between the different sections and the rationale behind all the different model comparisons and parameter fits require further motivation, especially for some important design choices.} The text has been reworded to better highlight the fact that
        \begin{itemize}
            \item I adopt a probabilistic perspective, where the goal is to study and come up with a model for the joint distribution MT, ID$_e$, parametrized by strategy.
            \item I study this at a local (within condition) and global (across conditions) levels
            \item The objective is dual: to better understand the speed-accuracy tradeoff (explain), but also to be able to generate realistic pointing data (predict)
        \end{itemize}
    The focus is now less on the methods (copulas and EMG), and more on the logical flow and about what we learn.
    \item  \textit{The paper lacks a general motivation and is rather vague in the conclusions that can be drawn, which limits its immediate value to the HCI community.} The motivation is to better explain and predict pointing, as should now appear more clearly. The two major results are now better highlighted in the discussion
        \begin{itemize}
            \item \textbf{Predicting pointing data from nominal D, W and strategy values, with parametric models.} I introduce several models that can be used to generate realistic pointing data: a model that predicts ID$_e$ values based on strategy and ID, a bivariate Gaussian model that can be used to generate block summaries, and three models that can be used to generate pointing data, for given strategies and ID levels. The three models are also compared on metrics that are directly used by Fitts' law researchers.
            \item \textbf{A distinction between the strong global speed-accuracy tradeoff and the weak local speed-accuracy tradeoff} While this knowledge is not leveraged in this work to produce direct implications for design, it contributes to a better understanding of pointing.
        \end{itemize}
    \item \textit{I consider Sections 5 and 10 particularly important, as they provide an overview of the study design and results. However, it is not entirely clear to me how the two approaches from Section 5 relate to the three models ultimately proposed. For example, why are bivariate Gaussian models not mentioned in the beginning, and how do they relate to copulas?} Model 1 relates to the approach that directly models the joint (with copulas), while Models 2 and 3 relate to the approach based on the conditional distribution. The bivariate Gaussian model is now introduced in the ``overview'' section. The bivariate Gaussian is a model for the joint distribution between ID$_e$ and the \textbf{average} movement time, whereas the copula models describe the joint distribution between ID$_e$ and movement time.
    \item \textit{Most of the statistical models considered aim to derive correlations between MT and ID$_e$, as well as between ID$_e$ and one or multiple of the independent parameters from the protocol (D, W, ID, s). It would be nice to see arguments for these design choices. In particular, the main assumption that MT only depends on ID$_e$, but not on D, W, or s (Eq.(19)) requires further explanations.}  We could indeed learn the joint distribution of MT and ID$_e$ conditioned directly on D, W, ID, s. However, one would need to come up with a model that accounts for all these term, which would likely become something overly complex, prone to overfitting, and difficult to understand. Since ID$_e$ is a function of D,W and strategy (via the error rate, see explanation in 2.2.2), it seems reasonable to study separately the effect of D, W and s on ID$_e$ on one hand, and how this value of ID$_e$ affects MT on the other hand. Technically, the assumption is that MT $\perp$ (D, W, s) $\mid$ ID$_e$, which means that if one knows ID$_e$, then adding D, W and s adds no information to MT, or in other terms, blocks that have the same ID$_e$ are statistically similar when it comes to predicting MT values. This remains a modeling choice, chosen because it favors simplicity over adding ``second order'' terms.  This is now more explicitly mentioned in the paper: ``\textit{This reduction avoids considering too many variables and we justify it by the fact that ID$_e$ already summarizes task geometry and strategy \ie the remaining effect of s, D and W on MT beyond those due to ID$_e$ are likely second order.}''
    
    
    
    
    \item \textit{Also, why is ID$_e$ instead of ID used, especially in the EMG model? The effective index ID$_e$ is a dependent variable not known in advance, which renders the conditional distribution p(MT|ID$_e$) a purely observational, posterior quantity, whereas EMG (and other forms of nominal Fitts' Law) can be used to predict the expected movement time for different ID values. } This is precisely why the model of ID$_e$ conditional on D, W and s is introduced, all three latter quantities being known in advance, before movement is being recorded. Together with the ID$_e$ model, p(MT$\mid$ID$_e$) can be used to predict full datasets (in two steps: first predict ID$_e$, then feed ID$_e$ to the conditional distribution). This is now better highlighted.    
    EMG has been used before on both ID and ID$_e$: the effective index tends to shift low ID to the right and high ID to the left (which explains the next point), which increases a bit the density of points, but it does not fundamentally change the shape of the distributions. 
    \item \textit{Notably, "the nominal version of Fitts' law is strongly preferred over the effective version" for the third dataset (Section 8.2); why is a similar comparison not conducted for the first dataset (Section 6) as well?} It is a well known fact that in most Fitts' law datasets the nominal version leads to higher $R^2$ than the effective version. This is usually due to having the same variability of data on the y-axis but in a smaller range on the x-axis due to the ID$_e$ correction. I have decided to remove this comparison in the work since it is largely off topic, it was here as some form of preliminary result, to make sure the dataset is consistent with what is expected.
    \item \textit{I enjoyed the comparison between local (condition-dependent) and global (condition-independent) effects for the considered statistical methods. Even though hardly any local correlations between movement time and accuracy were found, these findings contribute to a better understanding of the speed-accuracy trade-off across study types/conditions. I feel that the interpretation of these findings should be placed more prominently, \eg, in a discussion section on the different protocols and the variance between individual users and/or prescribed strategies.} This is now adressed in the new discussion subsection on the local and global tradeoff, and the local-global tradeoff is also emphasized from the introduction.
    \item \textit{Given the large number of statistical models tested, most models don't show a large or significant parameter correlation.} Correlation is a property of the data and is independent of the model used to represent said data. I think here the reviewer might have misunderstood what model evidence ratio (MER) used to rank copulas represents. MER has no bearing on correlation, it just tells us given observed data, what are the odds a particular model is better than the current best one. The very low MERs for many copulas just indicate that almost without a doubt, these copulas are inferior models. However, they might still indicate strong or weak correlations between ID$_e$ and MT. The correlations are in line with the local/global SAT: correlations within conditions are usually weak, correlations between conditions are usually strong.
    \item \textit{It is unclear which model fitting methods the authors suggest to apply} The goal of model fitting can be either to better understand the dependencies between ID$_e$ and MT (explain), or to leverage this to better simulate pointing data (predict). From a prediction perspective, the EMG models perform better. From an explanation perspective, I prefer the copulas. 
    \item \textit{and which benefits the use of copulas might offer to HCI (given the many low correlations, and that the best fit usually is observed with t-/Gaussian copulas, which mainly predict a correlation between MT and ID$_e$ for medium values)} I now discuss the applicability of copulas in 11.3, and how and when they may be found useful. %Both the Gaussian and t copulas model symmetric dependencies. The t-copula captures tail dependence: extreme low or high values tend to occur together more often. As the degrees of freedom parameter ($\nu$) of the t-copula increases, causing it to approach the Gaussian copula, the tail dependence weakens, eventually disappearing when the t-copula becomes Gaussian.
    \item \textit{Even more confusingly, although low tail dependencies are found for some conditions and protocols (an interesting result), copulas are generally criticised because "the lower tail of the MT distribution should associate strongly with the entire ID$_e$ distribution" (Section 10.1), an expectation never discussed before.} This expectation comes from the ``minimum'' Fitts' law. The point of this sentence was to explain before the evaluation what could be the failure points of copulas, so that we have some expectation about the upcoming evaluation. Both points have now been made more explicit.
    \item \textit{It seems the paper would have benefitted from a simpler and more coherent evaluation, e.g., by collecting own participant data. As mentioned in the paper, there is a range of potentially confounding variables, including different input modalities used in the three separate datasets. Moreover, the parameters D, W, ID were not sampled independently in the study designs, which further complicates a fair comparison between identified parameters and correlation coefficients (and might have contributed to the lack of more significant results).} I disagree with this assessment, since the results are quite stable across different input modalities and protocols, which shows the robustness of the analysis and makes it more likely the results generalize. D, W are fully crossed in 3 out of the 4 experiments, as is the case in the vast majority of the Fitts' law experiments that are conducted in HCI. Yes this creates a collinearity between D, W and ID, but this has always been the case in Fitts' law studies. In all cases, I show that the effects of D and W can be neglected. This is also confirmed by the fourth experiment which uses a shape x scale design. Even if the coefficients are significant (in part due to the big sample size for that experiment), their actual values are very low. 
    \item \textit{The paper seems to pursue a purely statistical approach, with focus on a posteriori dependencies between measured movement times and effective IDs/accuracy. This definitely has value to the community, but should be further emphasised} This is now more explicitly said in the introduction and abstract.
    \item \textit{A further discussion on the expected applicability of the proposed statistical methods (e.g. copulas) to other laws such as 2/3 Power Law, Schmidt's Law, or Steering Law would further strengthen the paper.} I have included more discussion on the applicability of copulas to HCI. Unfortunately, while these are definitely interesting avenues, I can't currently make any strong links with the suggested laws, so I prefer not to mention them explicitly.
    \item \textit{It is unclear to which extent which of the identified parameters are meant to be reused or even act as default parameters for generative models. Notably, most parameters are primarily mentioned in the figure titles and legends, but then prominently referenced in the models introduced in Section 10. Albeit the authors are confident that "the parameters of the YKORM dataset [...] are likely applicable to many HCI contexts" and find they "could be reused directly by other researchers", their general applicability requires further discussion. Also, the steps required to apply the copula model in practice (Section 10.1 (1)) should be described in more detail (see further comments below).} The parameters have been estimated from the available datasets. As far as the datasets have been collected in conditions similar to the situation of interest, then the parameters would apply. Fitting the copula model is as simple as calling the relevant function in a library that supports copulas.
    \item \textit{The related work section covers the most important literature, but could be further extended. A more general perspective on formalisations of the SAT beyond the prevailing Fitts' Law-style mapping from amplitude and accuracy to movement time would be beneficial (e.g. Schmidt's Law [1]). Also, Section 2.3.1 should cite works that fit the aggregated (mean) version of Fitts' Law.} In line with comments by other reviewers, the related work section has been shortened. I did add mentions to other classical accounts of the tradeoff, like Woodworth, Schmidt, Elliott. Most works fit the aggregated version of Fitts' law, it is very rare to see the actual pointing data in a paper.
    \item \textit{Recent HCI models of perceptuomotor movements such as the multiple process model [2] should be mentioned (optionally, also traditional predictive models such as Meyer's optimized submovement model [3]). In addition, Section 2.5 should mention that there are simulation (e.g. PPO-based) approaches that can predict movement times for a specific study design (e.g., a dwell time protocol) and were shown to account for Fitts' Law; see e.g. [4, 5].} References have been included and discussed (except Meyer's).
    \item \textit{The paper has medium originality. Although copulas (and some other combinations of statistical models proposed) might not have been used in HCI before, it remains unclear to which extent they drive forward Fitts' Law-related research. At least, the authors should refrain from claiming copulas a "novel statistical tool" in the abstract, given that they do acknowledge that "the tools to work with copulas remain somewhat underdeveloped, which may add barriers for some researchers" ('statistical method' might be more appropriate).} I think papers that model pointing are never considered original by reviewers. I suppose it is seen as a "done deal" in HCI. I personnally believe that jointly considering MT and ID$_e$ distributions is original for two reasons: modeling distributions is rare in HCI, and considering the dependencies between two full distributions is even rarer. The proposed approach is also a nice sweet spot, in between models like Fitts' law which may be too simple, and simulation models that are convoluted and costly to train, and for which we have no guarantee of transfering to other conditions. I have moved the focus away from copulas in the abstract and introduction.
    \item \textit{Section 2.3 would benefit from further interpretations of the EMG assumptions and parameters (given that it is fundamental to this paper)} I have added more interpretations.
    \item \textit{Copula procedure (Appendix C.2): why is the sum rather than the average used to aggreagte likelihoods, if NaNs are counted as zeros anyways?}This boils down to how averages deal with NaNs. In practice, NaNs are not taken into account in the average of most methods \eg numpy.mean(). For example log-likelihoods of [30, Nan, 35] would average to 32.5 and sum to 65, whereas log-likelihoods of [32.5,32.5,32.5] would average to 32.5 but would sum to 97.5. The point here is that a failure to fit indicate a mismatch which should be penalized, which the average does not.
    \item \textit{- Section 5, "Dependence between..." paragraphs: a short explanation on why average MT is considered for the pure strategy protocol dataset, and MT for the other two datasets, would be helpful} The reason for considering average MT is now made explicit: it is to get an in-between model between the full distribution and Fitts' law, that is simple and interpretable. unaveraged MT are also considered in the analysis of the pure strategy protocol data.
    \item \textit{Section 7.2: ID$_e$ is a continuous variable, in contrast to the experimentally controlled ID; are ID$_e$ values binned when computing average MT "per block and ID$_e$"? (how?)} ID$_e$ can only be computed at the block level. It is obtained by considering all trials in a block, then computing sigma from these trials, and then applying the ID$_e$ formula. There is no binning involved.
    \item \textit{Section 10: examples require more detailed explanations of how to derive the copulas in practice (e.g., using Kendall's tau); in Model 1, the use of the unconditional EMG model for the MT marginal requires further explanations; in Model 3, the correction of lambda 1 should be better motivated} The copula parameters are estimated using maximum likelihood estimation. However, for those that don't master this technique, it is usually possible to estimate some copula's parameters from rank correlations. This has been made more explicit in 11.3. In Model 1, using copulas requires specifying marginal distributions, not conditionals, and so it is not possible to specify the EMG model with increasing variance and mean. The correction of lambda 1 is motivated in footnote 20.
    \item \textit{- Eq.(48): lambda 0 seems to be missing on the RHS?} Yes, added.
    \item \textit{- p.6,l.265: is sigma² the same as s in Eq.(10)?} Yes, but this does not appear anymore.
    \item \textit{using subindices to state with respect to which random variables the expectation is taken (with additional explanations introducing this notation) could further improve comprehensibility, especially for HCI readers with only basic mathematical background;} This is the ``standard way'' of presenting such results, to avoid lengthy notations. The laws of total expectation and variance have their own accessible Wikipedia page, with the same notations I'm using, so readers can easily look those up. I'm afraid adding the subindices will overload the notations, making it harder to understand.
    \item \textit{the proof, which is an essential part of the paper, could be further simplified by starting with a more intuitive interpretation of the "law of total variance", describing the contribution of its two components} The proof builds equally on the three ingredients (pearson correlation, laws of total expectation/variance), but I did add a footnote on the law of total variance to convey a more intuitive version of the law: ``\textit{This formula splits the total variance Y into the variance of group (X) means and the average residual within group (X) variance. It is conceptually similar to a one-way ANOVA, where the total sum of squares is partitioned into within and between group sum of squares.}''
    \item \textit{- Eq.(78) RHS lacks a square root applied to the entire fraction; the "-" in Eq.(80) should be "+"} Yes, modified.
    \item \textit{- p.8, footnote 9: "*strictly* convex/concave" functions have greater/smaller values} Yes, but I have removed this footnote altogether.
    \item \textit{- Fig.7 lacks a caption} caption added.
    \item \textit{- Section 7.4/Table 2/Fig.9 inconsistently use r and rho for (different) Pearson correlation coefficients} I use r when computing sample correlations, and rho when defining a parameter of the model. Fig 9's label on the right panel was indeed inconsistent, this is now fixed.
    \item \textit{- p.24,l.1206: should be strategy 3 that is highly significant, rather than strategy 2} Yes.
    \item \textit{(- Fig.20: why not show Gaussian instead of t copula instead (given that Gaussian has slightly higher model evidence ratios)?)} Because the Gaussian copula is a t-copula with $\nu \to \infty$, so the $\rho$ parameters are almost identical. I did remove this plot as well.
    \item \textit{- Fig.27: in caption, replace "rows" by "columns"} Yes, modified.
\end{enumerate}

\subsection*{Reviewer 2 comments}
\begin{enumerate}
    \item \textit{This starts with expectation management: From the abstract, I expected the paper to be about introducing copulas to HCI and using them to model the distribution of movement time in pointing tasks. While certainly true to some extent, in chapters 6-11 the focus shifts toward a collection of various methods (including copulas) and how they fit distributions to existing data sets, and a less definitive statement about how useful copulas are in this context compared to other methods. While the paper title fits very well, an abstract that emphasizes more that the paper is really about (fitting) various pointing models, and copulas are "just" one of them (line 93: "We also use copulas"), would make it clearer to the reader what to expect. } Yes I agree with this assessment and have rewritten the abstract entirely. I have also made the title more explicit.
    \item \textit{To put it bluntly, the paper sometimes reads like a collection of results (and the contribution list in the introduction seems to reflect that), with insightful gems of interpretation hidden in paragraphs, but without a clear structure and conclusion - particularly if one reads the paper for the first time. It feels like content was added at various sections at later stages, which, while more clear after having read the full paper, it is often confusing for a first-time reading. For example, in the contribution statement, the reader does not know (yet) about local and global dependencies (up until Section 5). Likewise, the EMG model has not been presented yet and moreover could be misinterpreted as having to do with electromyography. Throughout sections 6-8, I was wondering what is the benefits of EMG and Gaussian Bivariate models, i.e., "Is there a clear winner model for explaining local or global SAT?"} Again, I agree with this assessment. The abstract and introduction now hint to the local and global tradeoffs, and the distinction between explaining (first part of the paper) and predicting (second part of the paper) should also be more clear now.
    \item \textit{This brings me to the next point. Sometimes, I was missing a common thread and was wondering why I am reading this section} I have tried to write better ``bridges'' between each section, so that the logical flow of the paper is clearer, as well as provide more information so that the reader can expect the content.
    \item \textit{"Why is the WHo model important? Why are its axioms important? Do the proposed models also need to satisfy these?". Readers unfamiliar with the model might also wonder what are x and y in axiom 3. Similarly, for HCI people without a deeper mathematical background, in eq. (6) the derivative dy/dx could be explained (if the WHo model even needs to stay in such detail).} I agree the WHo model was unnecessarily detailed, I shortened it a lot.
    \item \textit{Similarly, after having read about the EMG model, I was wondering about the main takeaways, and why I am reading about the model in such detail - until I saw the EMG model applied later.} I now explain the EMG model from a higher level, and also manage expectations with respect to EMG in the abstract and introduction.
    \item \textit{Another example is Section 3.1, which I personally like, but at the same time I wondered, "Shouldn't the reader be already convinced by the intro, i.e., shouldn't these arguments be made in the intro already?".} I removed 3.1, I agree it came a bit late to motivate modeling distributions of MT. 
    \item \textit{Yet another example is Section 6.3.1, where it feels like results are presented but then immediately discarded ("there are only four distinct ID$_e$ levels, which we believe is insufficient to yield robust results", line 706).} I reworked the results section and this no longer appears.
    \item \textit{Similarly, I would wish for more motivation as to why, after the results from 6.3, one needs/wants to continue with 6.4.} The point of 6.4 is to have a model for IDe based only on nominal values p(IDe $\mid$ D, W,), so that the models that only build on the conditional p(MT$\mid$IDe) can be fully predictive. This is now said more explicitly in section 4.
    \item \textit{Moreover, the "revisiting" previous results using yet another data set in Section 9.3 appears too sudden and not well-motivated. Generating data in Section 10 also feels like a sudden topic switch.} I believe revisiting results with the PDD dataset is a strong point; Reviewer 1 also appreciated it. It allows dispelling some ``doubts'' that we could have regarding previous results (input techniques which were not purely standard mouse, no small W's, shape x scale design). I have made it more explicit what this re-analysis brings.
    \item \textit{ Generating data in Section 10 also feels like a sudden topic switch. } I now more explicitly state that there are two goals in this paper, using Shmueli's words:  to explain and to predict.
    \item \textit{Some parts of the paper feel rushed. While not all figures have to be perfect for review, it would still be nice if the paper felt more "publication-ready":} I have reworked some figures as well as the captions
    \item \textit{References to "subsubsections" or figures often miss a word in front of them} This is poor english grammar from my part, I have now added the needed word.
    \item \textit{- Section 12.3 feels incomplete.} Future work is now integrated in the conclusion.
    \item \textit{Overall, while each paragraph is clear, the paper as a whole is hard to digest in its current form because of the reasons above. I therefore recommend a major overhaul regarding the structure. Motivating things or leaving them out would make the paper more concise. I sincerely hope this review helps the author to improve their draft in these regards.} I have performed a major overhaul as I agree with this assessment. I have motivated the approach better and also left things out to focus on what is important.
    \item \textit{- line 265: Is lambda the same as in eq. 6?} No, but both have been removed.
    \item \textit{- eq. 11: What is lambda here?} removed
    \item \textit{- In eq. 14 and 15, ID$_e$ formally is both a random variable and a realization of one. I am not sure if this might cause confusion, and whether it should rather read E[MT$\mid$ID$_e=$ some specific value]. } The notation allows ID$_e$ to be both a random variable or its realization. Typically, one would use lowercase for the realization say X for the random variable and X=x for the realization. Here the equation still makes sense for ID$_e$ as random variable, and writing ID$_e$ = id$_e$ would be particularly ugly I think.
    \item \textit{- section 2.5.2: The distinction between simulation models and optimal control problems seems a bit arbitrary. Optimal control models that predict trajectories via Deep Reinforcement Learning might take hours or days to train, similar to reference [50]. Both might have real-time inference. I suggest to split between models that predict whole trajectories and models that predict MT (and its variability) directly. } Optimal control models as described by E. Todorov and M. Jordan occupy a special place in motor control, being particularly popular, which is why I alot them a distinct paragraph.
    On optimal control and the link with deep RL, this was actually pointed out by Reviewer 3 who prefered separating the two concepts. The two concepts are quite close, and some fields will consider optimal control problem with a more or less large scope. The separation boils down to the fact that OFC is a standard method in motor control that has closed form solutions but has not demonstrated realistic movement time distributions, while the simulators are much less standard, can be quite contrived, take up a huge amount of ressources for training, and can likely generate movement time distributions but we have no way of knowing if they are realistic without actually using the simulation model.
    \item \textit{- line 398: contradicts line 241, where the correlation is just r, not r²} Yes, edited.
    \item \textit{- line 417: What happens if k and theta *both* increase with ID$_e$? } If they both increase linearly with IDe, then we get a quadratic increase for the product $k\theta$, which does not fulfill the condition anymore. They could both increase with a square root of IDe. This type of reasoning shows what assumptions one can make on model parameters to have distribution of movement times that still satisfy Fitts' law. I have added this example, as I felt it illustrated well how Proposition 3.1 can be used in practice.
    \item \textit{- line 620: I wonder whether this assumption should be communicated earlier, not in a formula, but in text. Also, does this assumption apply to approach (2)? } No this does not apply to approach 2 since there we directly learn the joint distribution.  I have not communicated it earlier but did explain a bit more in the text what this assumptions entails.
    \item \textit{- line 648: What is the motivation of using MT in studies 6 and 8 but the averaged MT in study 7? } The motivation for the bivariate gaussian model for average MT and IDe is now better explained.
    \item \textit{- line 682: "a difference of 10 is usually considered good" - Is there a reference for that? } It's a classical rule of thumb by Anderson and Burnham (0-2 little evidence of a diff, 4-7 moderate evidence, 7-10 strong evidence, $>$ 10 very strong evidence). I have added that reference explicitly when this number is given (the book was already referenced in the paper).
    \item \textit{- line 693: What is the interpretation of R=1? } R=1 means 1:1 odds, or equal evidence that model A is better than B. Indeed the previous phrasing for describing R was not entirely correct, it is now changed.
    \item \textit{- line 863: A comparison is made and immediately after, footnote 14 tells that these are "not necessarily meaningful". This is confusing. } Yes, but at the same time it is the reality. Comparing pixels is not necessarily meaningful, and we authors don't always report the information to allow for meaningful comparisons (mostly DPI and gain settings of the mouse).
    \item \textit{- line 930: Without knowing reference [71], what is I$_u$? }It is the index of utilization as defined by Zhai et al 2004. But I have edited this section and it does not appear anymore.
    \item \textit{- Section 7: The structure of this section is different from Section 6. I would have expected a similar structure for comparison. } The pure strategy protocol does not have D and W, but only strategy as independent variable, so the analysis is different from the traditional Fitts protocol.
    \item \textit{- line 996: Theorem 3.1 -> Proposition 3.1} Yes
    \item \textit{- Figure 8: The different colors are hard to distinguish. This is true for a few figures (e.g., fig 14) and does not help accessibility. } Rather unfortunately, this is actually the colorblind palette of seaborn \verb#colorblind_palette =# \verb#seaborn.color_palette(palette="colorblind")# which I chose for accessibility reasons. I will gladly use another palette if there is a more recommended one in a revision.
    \item \textit{- line 1329: Would it make more sense to delete data with distractors to closer align to the other studies? } There were distractors in all conditions, just their density was varied. As shown in \textbf{Blanch, Renaud, and Michael Ortega. "Benchmarking pointing techniques with distractors: adding a density factor to Fitts' pointing paradigm." Proceedings of the SIGCHI Conference on Human Factors in Computing Systems. 2011.}, distractor density barely affects regular pointing. This is now explained as well in the paper.
    \item \textit{- line 1353: What is "one unit of strategy"? Giving the strategies numbers between [-1,1] seems arbitrary. } I think that is one of the force of the approach: mapping strategy to a number in the [-1, 1] interval, where -1 is the maximum speed strategy completely disregarding accuracy, and 1 the opposite. Of course [-1, 1] is arbitrary ([0,1] would do just as well), but this has little bearing on the remainder, it just changes the estimated parameter values. This mapping indeed leads to a useful representation, since it leads to a simple linear model for the locations of the ellipses, consistent with Fitts' law. I have highlighted this more.
    \item \textit{- Section 9.1: This is a nice description of what was done, but I found the motivation lacking. } I have made the synthesis sharper.
    \item \textit{- line 1809: "we estimated ..." <-- where?} Here I am reporting a new result.
    \item \textit{- line 1850: Please recall where the alpha and beta values are documented. } Yes, edited.
    \item \textit{- line 2171: "which for all models was good": Was it? It can generate negative values. What constitutes "good?"} This part was rewritten.
    \item \textit{- Figure 26: What do we learn qualitatively from these? Especially since some models have a negative, others a positive slope.} This part was also rewritten. The point is to verify the alignement with the gaussian bivariate model.
    \item \textit{- eq. 60 is empty. } removed
    \item \textit{- line 2273: Recall where "local" and "global" dependencies were "defined". } now defined in abstract, introduction, outline
    \item \textit{- References: Please update references to arXiv preprints to their respective peer-reviewed publication, if available (e.g., reference 60). } I only found this one, which I have fixed
    \item \textit{- Appendix B: I have read the proof and I believe it is correct. I have some minor remarks. Do X and Y have to be random variables on the same probability space for eq. 62 and 64 to work? If yes, recalling that it is about *joint* distributions might help. } Yes technically, but in my opinion it is way to technical for an HCI audience to start giving this level of precision. You can always construct a common probability space for two independent distributions anyways, and then the formulas simplify straightforwardly. I prefer not to add this precision.
    \item \textit{- Appendix B line 2550: "always larger" might need to be switched to "always at least as large"} Yes
    \item \textit{- Appendix B line 2556: A reference to "a linear conditional expectation" could be helpful. } It is defined in Eq. 49
    \item \textit{- Appendix B eq. 78: Y instead of y. } yes
    \item \textit{- line 2622/eq. 81: Is k "the number of parameters" from line 2620?} yes, fixed
\end{enumerate}

\subsection*{Reviewer 3 comments}
\begin{enumerate}
    \item \textit{The main flaws are the focus on effective ID, which makes models dependent on user data,} I have adressed this point in the previous comments to reviewers, but as a reminder this is the whole point of the $p(ID_e|D,W,s)$ model, which makes the model able to generate the full joint distribution.
    \item \textit{lack of a proper discussion of the used method} Copulas are now discussed further in 11.3
    \item \textit{contrasting previous findings} there is limited work on analyzing and predicting the joint distribution, but I have compared the models to Fitts' law whenever possible. The IDe model is compared to Zhai et al.'s 2004 work.
    \item \textit{missing direct evaluations against previous methods or real user data} There are to my knowledge no previous methods that I can compare to. If the reviewer points out a parametric alternative/baseline which accounts for strategies and predicts entire distributions that is reasonably implementable, I would be happy to use it as comparison. As for ``real'' user data: The data I used is not fake, simulated or whatsoever. It is just data from existing controlled experiments. That does not make it less valid, especially considering running a new Fitts-like pointing study would probably add little, except if very specific conditions were to be tested, which I don't think there is a need for here.
    \item \textit{or, the application of the presented approach to a real design problem.} It is true that there are no concrete \textit{implications for design.} However, the work offers better understanding of pointing and the speed-accuracy tradeoff, and has a few concrete contributions that directly help those that are interested in modeling pointing data: a theoretical result that guides the selection of distribution of movement times (subsection 3.1) and 5 new models: one for predicting ID$_e$ from nominal conditions, one for predicting the joint average MT, IDe distribution, and three for predicting the joint MT, IDe distribution. These models could be used for a variety of things in HCI, including approaching a real design problem, typically any optimal design problem where part of the task completion time is due to pointing. 
    \item \textit{Using copulas in the first place is weakly motivated and it is not explicitly supported by the results} I have now better motivated copulas, as a way to investigate the dependence between IDe and MT without any assumptions. I'm not sure to understand the criticism that copulas are not supported by the results, since I am only performing model comparisons. Thus, the results show that some copulas are better than others at modeling the data, not that copulas are a good or bad fit.
    \item \textit{The paper is well written and contains excessive statistical comparisons of the proposed probabilistic models.} I'm not sure how to interpret this, but the number of statistical comparisons has been reduced.
    \item \textit{The paper’s structure is slightly confusing. Many different ideas and models are introduced, which are then evaluated on different existing datasets which makes it difficult to follow the flow of the paper. This is also reflected by the variety of main contributions the paper claims to have. } This is in line with the previous reviewer's assessment and I agree, and have tackled this by a major rewrite, including the structure of the paper.
    \item \textit{The paper unfortunately falls into the “yet another Fitts’ Law paper” category. The paper presents new models with new parameters, but since the evaluation is flawed (see below), it is unclear whether they will find application in further work. The motivation of presenting another model for movement time when models for complete movement trajectories -- which consider biomechanical constraints and allow to analyse human movements in a much broader way than summary statistics -- already exist is unclear. } I don't agree with this assessment. While Fitts' law aims at modeling average movement times and aims to make the variability of MT dissapear, I explicitly model this variability. There are no summary statistics, except in the Gaussian bivariate model. Compared with biomechanical models, there are advantages and disadvantages. For example when doing optimal design, setting up an entire biomechanical model, which has to be retrained for each particular condition seems like a big overkill to me. Also, it is not clear how strategy is incorporated in the biomechanical models I've seen, but maybe there is a nice way to do it. Also, the generated MT distribution of biomechanical models have to my knowledge not been validated, they are usually evaluated on summary statistics \eg showing they can replicate Fitts' law. In general, there is no single best tool, and it may be advantageous to have other tools than biomechanical models. 
    \item \textit{Although the introduction of copulas as a novel statistical tool into HCI and the analysis of pointing behaviour is interesting, the other contributions are of low originality. No new user studies were performed.} I think the originality lies in the analysis, not in the data collection / experimental aspect.
    \item \textit{The paper mentions different related models and covers the area well. However, there is no direct or explicit comparison of the proposed models against previous models. Also, optimal control and reinforcement learning models are mixed up.} I don't have access to the simulator models, and while I have myself implemented several optimal control models and could actually test some predictions, at least two things are unclear: 1) how to vary strategy (potentially by changing the weights allocated to the various components of the cost function) 2) how to define the end of movements (either use a fixed horizon, in which movement time needs to be supplied in advance, or use a receding/infinite horizon but then what is the stopping criterion?). To my knowledge there is no paper that discussed how to do this properly, so this would open an entire new can of worms. It is also known that optimal control models fail at properly capturing the magnitude of variability of human movements, both in space and in time, see eg. \textbf{Berret, Bastien, et al. "Stochastic optimal feedforward-feedback control determines timing and variability of arm movements with or without vision." PLOS Computational Biology 17.6 (2021): e1009047.}. So it is not easy to find a baseline to compare to. This is why I compare on the methods in how accurate they reproduce relevant features of the dataset in 10.1 and 10.2.
    \item \textit{I belief that the paper would benefit strongly from focussing on one aspect at a time, or even being split into different papers looking at copulas on one hand, and including strategy in the previous EMG model on the other hand} I put less emphasis on copulas in this revision, and removed some details or analyses that would detract from the main topic of study.
    \item \textit{Major parts of the work are based on the concept of effective ID $ID_e$, which is a combination of the classical ID of Fitts’ Law with endpoint variance. As such, for a given pointing task, it can only be computed a posteriori. This renders models that are based on given $ID_e$ unusable for movement time prediction. } As adressed in a few other points before, this is not the case since a model for IDe conditional on D, W and s is also provided.
    \item \textit{The model comparison in section 11 uses $ID_e$ values directly obtained from user data.} This was done to reduce noise for the analysis, since the most interesting and difficult is predicting movement time distributions. 
    \item \textit{The author themselves assumes that “the movement time distribution only depends on W, D and s through $ID_e$” (ll. 619, Eq. (19)). If this is the case, and the models know the exact $ID_e$, then it is no surprise that the models can recreate movement times well and that metrics are close to those of the real data} I have already discussed this point in other answers, but I want to stress that even knowing IDe, it still does not give one the distribution of MT just like that. For example, with Fitts' law, one gets MT = a + b IDe + Gaussian noise added to it. Clearly this would give data that would not fit at all with the empirical datasets.
    \item \textit{This demonstrate the general problem of using $ID_e$ in the first place. Even if a model that combines MT and $ID_e$ is created using the data of one task, it is unclear to me how and if this model can be applied to a different pointing task.} by having a model of IDe as a function of D, W and s. Actually, the strong part about this, is that if the task model changes \eg one wants to account for distractors, say $d$, only the IDe model has to be relearned $p(ID_e \mid D, W, s, d)$, while the MT $\mid$ ID$_e$ model could remain unchanged.
    \item \textit{In ll. 877-920, the author suggests that D and W have the main effect on MT, and that the “actual differences in execution are likely due to inherent variability of movements rather than planning and do not adhere to a local speed-accuracy trade-off”. This needs to be elaborated more. What is meant with “inherent variability of movements”? There are many causes for variability in human movement, such as constant and signal-dependent motor noise [A, B], or noise in the cognitive system [C]} It would be hard to elaborate more on this observation I think. I happens that within conditions there is little to no speed-accuracy tradeoff. Is this surprising? To me, yes. Can I confidently explain this? No. One hypothesis I make is that perhaps the variability of movement times is just so large that it masks the tradeoff that may occur. However, that suggest this variability is due to processes not linked to strategy, and likely indeed include constant and signal-dependent motor noise. The other hypothesis is that there is no tradeoff, since once D and W is "selected", all movements produced are equivalent. These two hypotheses can be true at the same time ie, there is no tradeoff, but noise is also quite high.
    \item \textit{What is the benefit of using $ID_e$ then?} I don't understand this point. The benefit of using IDe in general is that it provides a condensed description of D, W and strategy, all in one index. This is convenient for many reasons, and is the reason why Fitts' law is used in the first place.
    \item \textit{ Also, in almost all pointing tasks, accuracy is imposed by W (e.g. clicking a button). If all resulting “movements” (this should be movement times in my opinion) are then equivalent, why do we need these models?} First, it should be noted that the models can be used for a fixed strategy (just set eg. $s=0$ for a balanced strategy). Second, the strategy of users is directly linked to the cost of errors. Higher cost of errors lead to more careful pointing, see e.g. \textbf{Bertucco, Matteo, Nasir H. Bhanpuri, and Terence D. Sanger. "Perceived cost and intrinsic motor variability modulate the speed-accuracy trade-off." PloS one 10.10 (2015): e0139988.}, so realistic pointing tasks are likely similar to a Fitts-with-strategy protocol. It should also be noted that the second study which uses no target width is just one study among the 4, and it is used in part to more gradually introduce the analysis of the third study.
    \item \textit{There are no direct evaluations of the presented models against existing models or the user data itself (except for section 11). This makes it difficult to understand how well these models actually perform. The presented statistical methods are model comparisons, based on the model evidence ratio.} The first part of the paper deals with model comparisons, because I am assessing with copula family fits best, to get an idea of what the variability actually looks like. In the second part, centered around predicting data, I do use comparisons with the user data. This is in line with model evaluation/comparison practices (see e.g. \textbf{Shmueli, Galit. "To explain or to predict?." Statistical science (2010): 289-310.}). I have made it more explicit in the introduction that the first paper is about explaining/measuring, while the second part is about prediction.
    \item \textit{The final models presented, however, can lead to negative movement time for low IDs (ll. 2145), which obviously does not make sense. This issue needs to be discussed and possible solutions could be presented (such as using a log-normal distributions). } I have recognized this caveat indeed for the EMG based models. This problem is in practice not very important, since the number of negative movement times is low and occurs for ID values below 2, which in practice rarely occur. For example, Fig 9 shows that a single block has IDe below 2.5. Optimal design scenarios are usually focused on smaller targets, as this is where there is time to be gained. I'm not sure how the log-normal could be used: In practice, we would need the shifted version, which may produce negative values as well. However, a simple solution is to post-process the dataset and remove the few pointing times that are below, say, 200 ms. While it may seem inelegant, it is actually very affective and barely changes the features of the datasets. I have now included this extra step in the paper.
    \item \textit{Another point to discuss is that the evaluations are based on human data from existing data sets with slightly different tasks and different pointing techniques (e.g. gesture, bubble cursor, classic mouse pointing). On one hand, this shows potential that the models generalise well. On the other hand, however, this makes it hard to understand where problems in the fitting process or evaluation stem from. For example, a pointing task for targets without explicit width is very rarely seen in practice, and is quite different from pointing with a given boundary on accuracy (e.g., clicking a button). Conducting a new study with focus on different strategies would be very valuable for the work. } I see this as a strong point of the paper, which shows that the analysis is robust across various protocols and input modalities. The experiment suggested here is covered by the third study (Fitts with strategy).
    \item \textit{The copula fitting process seems overly complicated and is difficult to comprehend. } Copula fitting is based on maximum likelihood and is implemented in copula libraries. Some libraries have more families of copulas available than others.
    \item \textit{Section 9.1 describes parts of the process and discusses some results. I am uncertain what to think about the introduced bootstrap method for the rotated Gumbel, which is argued as followed: “at the time of writing the conditional form of rotated copulas was not available in the library we used”. Why was no different library used? Could this not be implemented? If not, how easy is it to actually recreate the presented methods? I appreciate that the code is provided, however, I am uncertain about the amount of effort needed to fit the models to a different dataset.} The boostrap method for the rotated Gumbel is a parameter/model recovery analysis: if the data is generated by X model with Y parameters, can I actually identify X and Y. This is a standard problem in modeling. Here the point was to verify whether the discrepencies between experiments could be due to some copulas being hard to identify or not. As now discussed in 11.3, there are a few libraries for copulas, but some provide many more copula families than others. The conditional copula is a technical detail that is not very important: the point is that the copula generates the joint distribution by default. If one provides values for one variable (like is needed for generating block data, where all values of MT within the block are associated to the same IDe level), it can still predict the other variable, but for this the conditional form of the copula has to be used. Unfortunately, conditional copulas can not be computed for rotated copulas in the R library I used. I circumvented this issue by a resampling technique. The point is minor since it gives results that are qualitatively identical to the first case, where the copula could be used normally.
    \item \textit{The related work section is confusing and mixed with background for the presented methods. It would be recommended to make a clearer cut between roughly related work, related work that is very close to the presented one, and the background and models that are actually used in the manuscript (e.g., the Exponentially Modified Gaussian). For instance, I do not understand why the WHo model is explained in such detail but then not used for comparison. } The related work has been completely reworked, with namely WHo much shorter and EMG better explained.
    \item \textit{The citation of [20] in section 2.5.1 probably means [21], since [20] does not have 36 pages. The work in [20], which uses a reinforcement learning approach and not classic optimal control, shows that the resulting movements are in line with Fitts’ law. There is also no requirement for these movements to have a fixed duration (which is also the case for the model presented in [45]).} Yes indeed, I mixed up the two Fisher et al. references. This part has been rewritten. OC does not require fixing movement duration in receding and infinite horizons to generate movements, but some stopping condition has to be introduced to define the end of movements to define movement time and endpoint. Fisher et al. in the 2022 ToCHI paper computed endpoint standard deviation by an arbitray 0.97 seconds cutoff. Other stopping conditions can be considered. The points is that to my knowledge, no OFC model has been used to predict realistic movement time distribution, especially parametrized by strategy. This point has been clarified. The link between OC and RL has also been clarified.
    \item \textit{First, how comparable are strategy instructions, such as “Perform the task as rapidly as possible” [68] (which are used here as equivalent for strategy) between different tasks, studies, and even users? } Of course, there is a risk in transforming an ordinal instruction  to a numerical scale, eg, is the difference between speed emphasis (-1) and balanced (0) the same as between balanced (0) and accuracy emphasis (1). What matters is that with assumption, we actually obtain quite a simple linear model that is consistent with Fitts' law. So the transformation is maybe surprinsingly quite effective and useful. This is now emphasized.
    \item \textit{The author mentions that these instructions are sometimes applied relative to the previous condition the user has gone through, meaning that the order in which instructions are given have an influence on the resulting movements. } Properly counterbalancing conditions in an experiment where participants are asked to change strategy conditions is not straightforward. In most previous experiments, all blocks featuring the same strategy were always contiguous; this is done to try and have homogeneous participant behavior for the different strategies. Thus, blocks grouped by strategy were counterbalanced across participants. For example, if the experiment total 36 blocks, 12 for each strategy, then the 12 blocks belonging to the same strategy were always contiguous for all participants, although the order in which they would face the 3 groups of 12 would be different.    
    Some examples include: In the YKORM dataset all participants first performed the balanced blocks, the remaining two groups of blocks (speed or accuracy) were counterbalanced across participants. In \textbf{Zhang, Mingrui Ray, Shumin Zhai, and Jacob O. Wobbrock. "Text entry throughput: Towards unifying speed and accuracy in a single performance metric." Proceedings of the 2019 CHI conference on human factors in computing systems. 2019.}, all participants ended with the balanced block, and fast /extra fast were balanced with accurate, extra accurate. The GO dataset is an exception in that strategies were not contiguous, and changed at each block, but following the order (and in reverse) of the ordinal scale of strategy. The impact of improper counterbalancing is likely an inflated range of IDe values, and has likely litte bearing on the MT distribution itself.
    \item \textit{Additionally, as mentioned in Section 7.7, the size of the target can have an influence on the user strategy, e.g., very large targets can lead to open-loop or purely ballistic movements, which can also be seen in the PDD dataset [51] ID 2 condition (leading to $ID_e$ < 1 as mentioned in 9.3.1). The strategy instruction, as defined here, needs to be independent of the task width and distance. Otherwise, the same strategy instruction might have a vastly different effect for different targets. This should be discussed.} The YKORM dataset fully crosses ID and strategy, and the statistical models also do not assume dependency between strategy and (D,W). As determined from the 2 other datasets (PDD and JGP), the effect of D and W on IDe other than through the ratio D/W is minor, and can thus be neglected. 
    \item \textit{The paper is also not consistent with whether strategy plays a role or not, which is probably due to the different datasets used. Running a new study would also help mitigating this issue.} The paper's premise is that strategy, D and W influence IDe, which in turn influences the MT distribution. This is verified on all 4 datasets. There seems to be no need for a new experiment or study.
    \item \textit{There is no explicit Discussion section, which is surprising for such a long journal paper. The used methods and made assumptions need to be scientifically questioned, and possible risks and challenges need to be raised. Section 9 offers a few bits of this and sometimes interesting discussions can be found “between the lines”. } A discussion has now been added.
    \item \textit{Besides the other concerns raised here, differences to earlier findings need to be discussed, and why the author did not choose to address related issues. For example, whether the differences in line 861 stem from differences in widths used in the different studies or the choice of study design could have been addressed by running a new study. } This was already adressed by re-analyzing the PDD dataset.
    \item \textit{1./2. This is not novel. There have been many studies on ID and MT. The comments on the difference between average MT and distributions of MT in Section 3.1 are appreciated, but this is not a major contribution in my opinion. To make it one, it would require are thorough study of whether using the average has major drawbacks in real applications and how distributions can be effectively used in design of pointing tasks/techniques. Section 3.1 addresses this question, but only speculates on possible advantages. How would a distribution of MTs actually be applied in an example such as the one described in footnote 10?} The novelty is focusing on entire distributions and strategy. The use of average task completion times to generate optimal design is generalized in HCI. It is valid when all operations that lead to full task completion time are linear, which is not always true. It should also not be forgotten that the other aspect is including strategy. We know people use different strategies depending on the perceived cost of making an error. This is never taken into account, because there was previously no model. This could be investigated next, and is now mentioned as future work.
    \item \textit{3. This is an interesting bit of the paper. To introduce copulas to the field of HCI, however, a solid foundation needs to be build. Section 4 is way too short and not very comprehensive. How does using copulas differ from using multivariate normal distributions in praxis? There are some comments on this here and there in the paper.} The focus is less on copulas in the revised version. However, to answer the question, two Gaussian marginals coupled with a Gaussian copula are equivalent to a multivariate normal. However, a Gaussian copula can couple any two marginals.  Fig 1 (e) shows a scatterplot obtained with the Gaussian copula which is nothing like the one you would get with a bivariate Gaussian.
    \item \textit{4. These models are only useful if they can actually be applied to real pointing tasks and are better than previous models. Questions regarding the strategy parameters arise (see above).} There are no equivalent previous models that I know off.
    \item \textit{5. This model is interesting, since it is fairly simple and probably easy to apply. Nonetheless, it does not include copulas which makes me wonder if it fits well in the rest of the work. Also, the fact that such a simpler model is performing similarly as the copula models raises the question whether such more complicated statistical tools are actually required.} The bivariate model is applied to the average movement time, whereas copulas are applied to the whole movement time distribution. The reason for using the bivariate model is now better explained.
    \item \textit{7. One should keep in mind that, to actually be able make use of a model of pointing (e.g. to evaluate an interface), the ability to generate realistic human-like data is paramount. The term “realistic pointing data” here is an overclaim – it should be “realistic movement time”. And although the results look good, the evaluation has a major flaw. The models are using $ID_e$ acquired from user data. This renders the evaluation pointless in regards to real application -- if we need to run a user study to be able to use the model, why use the model in the first place? The models can only create such realistic data because they know parts of the data, namely, the end-point variance that is defining $ID_e$. The calculated variable, the movement time, is not independent of this measure. } As explained before in this document, a model to generate IDe is provided. Movement time can not be independent of IDe otherwise there would be no speed-accuracy tradeoff, since the latter dictates higher values of IDe are associated with higher values of MT.
    \item \textit{Why is the minimum in the Model evidence ratio figures sometime 10-4, 10-3, or 10-2? } Because sometimes evidence is stronger for some copulas/datasets than for others. For example, increasing sample size usually leads to stronger evidence. Model evidence ratios can go all the way down to $0$ asymptotically.
    \item \textit{-	Throughout the paper, whenever a latex label is referred to, the “in” is missing (e.g. l. 488 “in subsection”).} added
    \item \textit{-	Fig. 7 is missing the caption} added
    \item \textit{-	Fig. 9  the values do not match the values in the text, e.g. r=0.35 although r=0.44. Please explain/fix. Why is strategy -1 only excluded for r and not for the other parameters? This seems hacky.} Considering -1 as special is now explained in the text. Essentially, movements without accuracy requirements are usually considered quite differently from aimed movements in the motor control literature, so it makes sense to separate that strategy.
    \item \textit{-	Many of the figures and tables are too wide.  Tab 1, 7, Fig 1, 9, 15,19, 25, 26} This is done on purpose to increase the visilibity of the figures, and would be modified by the typesetting of the journal I believe.
	\item \textit{-	Fig 25.: Consider removing the y-label and ticks on all but the first one (this might also help for other figures). Consider moving the legend out of the plots and only plot is once. The chosen dots for the scatter plot make it hard so see anything (especially in the second plot). Consider using smaller dots without white border.} The point of these datasets is just to visualize how variable pointing data actually is, considering we are usually confronted to nice clean Fitts' law models applied to block averages, and how the shapes of the datasets are roughly equivalent. The real comparison is performed by the features in the titles. The legends can not be shared since they embed the OLS and EMG fits proper to each panel.
	\item \textit{-	l. 125 EMG is not introduced yet} modified
	\item \textit{-	l. 195 erf not defined (later erfc)} erf-1 is defined in the cited reference. 
	\item \textit{-	l. 680 AIC is not introduced yet. } is now introduced
	\item \textit{-	l. 681 “largely superior” by how much? “routinely above” too imprecise.} The EMG model on this dataset has already been validated, so I don't want to spend too much time on it. Nothing surprising will come out of it.
	\item \textit{-	6.3.1: Why did the fit for some copulas fail exactly? The footnote is very vague.} For example, some copulas can only fit positive rank correlations. Because copula parameters can usually be computed from rank correlations, the initialization of the maximum likelihood procedure typically works by estimating the rank correlations first, and then transforming them into the copula parameters, after which a local optimizer is applied. This can lead to invalid parameter values. All of this is coded within the library.
	\item \textit{-	In 6.3.1 the participant data is pooled, but in the appendix C.2 it says that data is not pooled.} Data is not pooled, but the copulas are combined to form a single model evidence ratio; This is now made more explicit, and I have used the proper word "combine".
	\item \textit{-	6.4: I do not understand the first sentence. ID is only defined by D and W. how can D and W have an effect on $ID_e$ that ID does not have? I think what is meant is that the individual effects are analysed.} Indeed the individual effects are analyzed on top of the ratio D/W. I have made this more explicit.
\end{enumerate}

\end{document}